 \section*{第 \ref{ch:b3:quantum-angular-momentum} 章——量子角动量}

\begin{solution}{exo:b3:quantum-angular-momentum:1}
(a) $[\hat y\hat p_z - \hat z\hat p_y,\ \hat z\hat p_x - \hat x\hat
p_z]$：只有含 $\hat z, \hat p_z$ 配对的项保留，
得 $\iu\hbar(\hat x\hat p_y - \hat y\hat p_x)$。(b) $[\hat
L^2, \hat L_z] = \sum_i[\hat L_i^2, \hat L_z] = \hat L_x[\hat L_x,
\hat L_z] + [\hat L_x, \hat L_z]\hat L_x + (x \to y)$：四项
两两相消。(c) $[\hat L_z, \hat x] = \iu\hbar\hat y$：
$\hat L_z$ 生成绕 $z$ 轴的转动——对算符如此，
对态亦然。(d) 两个分量同时具有确定值，
由对易关系可知第三个分量亦须有确定值，且三者皆为零：
只有 $l = 0$ 能做到。
\end{solution}

\begin{solution}{exo:b3:quantum-angular-momentum:2}
(a) $\hat L_z = \hbar\operatorname{diag}(1, 0, -1)$。(b) $\hat L_+ =
\hbar\sqrt2\,\begin{pmatrix}0&1&0\\0&0&1\\0&0&0
\end{pmatrix}$，$\hat L_-$ 为其转置。(c) $\hat L_x =
\dfrac{\hbar}{\sqrt2}\begin{pmatrix}0&1&0\\1&0&1\\0&1&0
\end{pmatrix}$：特征多项式为 $\lambda(\lambda^2 -
\hbar^2)$。(d) $L_x = +\hbar$ 的特征向量为 $\tfrac12(1,
\sqrt2, 1)$：概率依次为 $\tfrac14, \tfrac12, \tfrac14$，对应 $m =
+1, 0, -1$。
\end{solution}

\begin{solution}{exo:b3:quantum-angular-momentum:3}
(a) $I = \mu r_0^2 = \qty{1.46e-46}{kg.m^2}$；$B = \hbar^2/2I =
\qty{3.8e-23}{J} = \qty{0.24}{meV}$。(b) $2B/h = \qty{115}{GHz}$，
$\lambda = \qty{2.6}{mm}$。(c) \qty{231}{GHz} 和 \qty{346}{GHz}。
(d) 谱线间距给出 $B$，进而得 $I = \mu r_0^2$：即
相距数光年之遥的分子的键长，其精度约为谱线间距相对精度的
一半。
\end{solution}

\begin{solution}{exo:b3:quantum-angular-momentum:4}
(a) $m\hbar$，其中 $m = -2 \dots 2$。(b) $\cos\theta = 2/\sqrt6$：
$\theta = 35.3^\circ$。(c) 完全准直将使 $L_x = L_y
= 0$ 与 $L_z$ 同时具有确定值，此为对易关系所禁止，
除平凡的 $l = 0$ 情形外。(d) $l/\sqrt{l(l+1)} \to 1$：当
$l$ 很大时，锥面收拢至轴上，经典矢量
得以复现。
\end{solution}

\begin{solution}{exo:b3:quantum-angular-momentum:5}
(a) 因不依赖 $\varphi$，算符给出
$-\hbar^2(\sin\theta)^{-1}\partial_\theta(\sin\theta\,(-\sin\theta))
= 2\hbar^2\cos\theta$。(b) 对同样的计算，计入
$\eu^{\pm\iu\varphi}$ 因子：本征值为 $2\hbar^2$；$\hat L_z$
给出 $\pm\hbar$。(c) $\cos\theta$ 与 $\sin\theta\,\eu^{\pm\iu
\varphi}$ 在 $\vect r \to -\vect r$ 下变号（$\theta \to \pi
- \theta$，$\varphi \to \varphi + \pi$）：宇称为 $-1 = (-1)^1$。(d)
它们是厄米算符 $\hat L_z$ 的本征态，
本征值互不相同。
\end{solution}

\begin{solution}{exo:b3:quantum-angular-momentum:6}
(a) $2J + 1$ 个态共享该能级（简并）；玻尔兹曼
因子则随能量予以压低。(b) 令乘积取极大：$\dd/\dd J = 0$
得 $2 = (2J+1)^2B/k_{\text{B}}T$，即题中所给的 $J_{\max}$。(c) 在
\qty{10}{K} 时：$J_{\max} \approx 0.8$，故 $J = 1$ 占主导，
\qty{115}{GHz} 与 \qty{230}{GHz} 谱线明亮；在 \qty{300}{K} 时：
$J_{\max} \approx 7$。(d) $T \approx 2B(J_{\max} + \tfrac12)^2/
k_{\text{B}} \approx \qty{310}{K}$——即一种转动温度计。
\end{solution}

\begin{solution}{exo:b3:quantum-angular-momentum:7}
(a) 上能级分裂为五个，下能级分裂为三个，二者
间距均为 $\mu_{\text{B}}B$。(b) 间距相等时，
$h\nu = h\nu_0 + (\Delta m)\mu_{\text{B}}B$ 且 $\Delta m \in \{0,
\pm1\}$：只有三个位置。(c) $\mu_{\text{B}}/h = \qty{14}{GHz/T}$，故在 $B = \qty{2}{T}$
分裂为 \qty{28}{GHz}，即 $\Delta\lambda =
\lambda^2\Delta\nu/c \approx \qty{23}{pm}$——凭
塞曼 1896 年的光栅即可分辨。
(d) 电子自身的自旋，其反常因子 $g \approx
2$：这些“反常”图样在自旋被命名之前，
就已是自旋的指纹（见\cref{ch:b3:spin-two-level}）。
\end{solution}

\begin{solution}{exo:b3:quantum-angular-momentum:8}
在所给单位下 $U_{\text{eff}} = -2/x + l(l+1)/x^2$。(a) 对
$l = 1$：极小值位于 $x = 2$（$r = 2a_0$），深度为 $-E_{\text{I}}/2$。
(b) $U_{\text{eff}} > 0$ 在 $x < 1$ 时成立：一个玻尔半径之内，
势垒占优。(c) 在原点附近 $\psi \sim r^l$：只有 $l = 0$ 的态
在 $r = 0$ 处密度非零。(d) 俘获电子需要
电子密度\emph{恰在}原子核处：唯有 $s$ 电子，
不受离心势垒阻挡，因而被吞没。
\end{solution}

\begin{solution}{exo:b3:quantum-angular-momentum:9}
(a) $\hat L_x = (\hat L_+ + \hat L_-)/2$ 会改变 $m$：对角元
为零。(b) 由对称性，两个横向平方
相等，其和为 $\langle\hat L^2 - \hat L_z^2\rangle$。(c)
在 $m = l$ 时：$\Delta L_x\Delta L_y = l\hbar^2/2 = \tfrac\hbar2
\,l\hbar$：取等号——极端态的准直程度已达
代数所允许的极限。(d) 锥面：$L_z$ 有确定值，横向均值为零，
横向涨落相等——一个长度与投影均确定的矢量，
在方位角上均匀地“涂抹”开来。
\end{solution}

\begin{solution}{exo:b3:quantum-angular-momentum:10}
(a) 由 $E_J = BJ(J+1)$ 作能量核算：$J \to J + 1$ 增加
$2B(J+1)$；$J \to J - 1$ 减少 $2BJ$。(b) $\Delta J = 0$ 被
禁戒，两支分别始于 $\pm 2B$：没有谱线落在
纯振动频率上。(c) 两组间距为 $2B$ 的梳状谱
夹着一个空隙，强度先升至 $J_{\max}$ 再降低。(d)
谱带的宽度即两组梳状谱中被占据的范围，随
$\sqrt{T}$ 增长：正是这些热布居的边翼，
即饱和温室吸收带持续吸收之处。
\end{solution}

\begin{solution}{exo:b3:quantum-angular-momentum:11}
(a) $\|\hat J_\pm\ket{j,m}\|^2 = \bra{j,m}\hat J_\mp\hat J_\pm
\ket{j,m} = \hbar^2[j(j+1) - m(m\pm1)]$。(b) 若继续升阶越过 $m = j$，
则得 $j(j+1) - j(j+1) = 0$，再进一步则模方为负，
故此链必含被湮灭的最高态。(c)
最高态与最低态相差整数个单位步长：$2j \in
\mathbb N$，即 $j = 0, \tfrac12, 1, \tfrac32, \dots$。(d) 该代数
从不涉及波函数；只有要求
$\eu^{\iu m\varphi}$ 在圆周上为单值，才迫使
$m$ 为整数——此条件只约束轨道运动，
而使内禀（自旋）角动量可取半整数。
\end{solution}

\begin{solution}{exo:b3:quantum-angular-momentum:12}
(a) 每单位体积为 $n\hbar$。(b) $\omega = n\hbar V/(\tfrac12 Ma^2)
= 2n\hbar/\rho a^2 = 2 \times \num{8.5e28} \times \num{1.055e-34}/
(7900 \times \num{e-4}) \approx \qty{2.3e-5}{rad/s}$——单次
测量中不可见。(c) 与纤维的共振同步地反复换向，
累积约 $\sim Q$ 次冲击：$\omega \sim \qty{2e-3}{rad/s}$，
周期为秒量级的毫弧度摆动——借助镜子
与光束即可观测，正如爱因斯坦与德哈斯所观察到的。(d) 实测的
旋磁比为 $e/m_{\text{e}}$，是
轨道值 $e/2m_{\text{e}}$ 的两倍：铁的磁性并非来自轨道电流，而是
电子的\emph{自旋}，其 $g \approx 2$ 由后续各章说明。
\end{solution}

\begin{solution}{pb:b3:quantum-angular-momentum:1}
\textbf{1.} $E_J = BJ(J+1)$，简并度为 $2J + 1$。
\textbf{2.} $I = \mu r_0^2 = \qty{1.46e-46}{kg.m^2}$：$B =
\hbar^2/2I = \qty{3.8e-23}{J}$，即 $B/h = \qty{57.7}{GHz}$——
即实测常数。
\textbf{3.} $115$、$231$、\qty{346}{GHz}。
\textbf{4.} $E_{J+1} - E_J = 2B(J+1)$ 线性增长：相邻
谱线相差常数 $2B$。由间距得 $I$，进而
键长——以无线电完成的结构化学。
\textbf{5.} $\lambda = c/\nu = \qty{2.6}{mm}$。大气中的水蒸气
强烈吸收毫米波，因此选址于阿塔卡马、莫纳克亚
与南极——高、冷、干。
\textbf{6.} 光子是自旋为 $1$ 的粒子，携带 $\hbar$：
总角动量守恒要求分子的 $J$
改变一个单位。
\textbf{7.} 其电荷分布在每个转动角度下都保持对称：
无振荡偶极子，因而没有偶极谱线——
堪称完美的伪装。
\textbf{8.} $B_{\text{H}_2} = \hbar^2/2\mu r_0^2 \approx
\qty{7.6}{meV}$；对称分子的弱四极
跃迁需要 $\Delta J = 2$，代价为 $6B \approx
\qty{46}{meV}$——为 \qty{10}{K} 时 $k_{\text{B}}T$ 的五十倍：低温
氢气根本无法辐射。
\textbf{9.} $2B = \qty{0.48}{meV} < k_{\text{B}}T = \qty{0.86}{meV}$：
碰撞使前几级保持布居——阶梯
恰好在分子云所在之处发挥作用。
\textbf{10.} $n_1/n_0 = 3\,\eu^{-0.478/0.862} = 1.7$：发射
能级布居充足。
\textbf{11.} $J_{\max} \approx \sqrt{0.86/0.48} - 0.5 \approx 0.8$：
布居在 $J = 1$ 处达到峰值。
\textbf{12.} CO 提供光子；将其换算为气体总质量
依赖于假定的 CO 与 H$_2$ 丰度比——
明亮的信使，标定却要付出赎金。
\textbf{13.} $\Delta\nu = \qty{115}{MHz}$：$v = c\,\Delta\nu/\nu
\approx \qty{300}{km/s}$，正在退行（频率降低）——星系
的轨道速率。
\textbf{14.} $\pm\qty{2}{km/s}$ 跨越 \qty{1.5}{MHz}；热
运动在 \qty{10}{K} 时仅给出 $\sim\qty{40}{kHz}$：宽度来自
湍流而非温度——谱线宽度描绘分子云的内部
“气象”。
\textbf{15.} 在中心质量周围，$v \propto 1/\sqrt r$ 本应
下降（开普勒）；而实测曲线的平坦意味着包围的质量
随半径不断增长——不可见的物质包裹着发光的
盘面：暗物质。
\textbf{16.} 布居之比，即
气体的激发温度。
\textbf{17.} 在 \qty{57.6}{GHz} 处：宇宙膨胀使每个
在途中的波长都加倍——正是本书末章
所讨论的宇宙学红移。
\textbf{18.} 光子通量 $\to$ CO 光度（所需距离）
$\to$ CO 数量（激发模型）$\to$ H$_2$ 质量（
CO 与 H$_2$ 之比的“X 因子”：最薄弱的环节，在局部标定并
谨慎外推）。
\textbf{19.} 温度为 \qty{10}{K} 的核心不发出光学望远镜
所能看见的任何辐射；其转动谱线同时是它唯一的
温度计、测速仪和标尺。
\textbf{20.} 温度计在能级间距 $\sim
k_{\text{B}}T$ 处最为灵敏：间距为 \qty{5.5}{K} 的能级在
$5$--$50$ K 范围内强烈响应，而 \qty{2}{eV} 的光学能级则
完全无从布居。
\textbf{21.} 谱线饱和（光学厚），更糟的是 CO
在最致密、最冷的气体中凝固到尘埃颗粒上：
天文学家转而使用更稀有的同位素分子（$^{13}$CO、C$^{18}$O）
以及其他分子。
\textbf{22.} 每种分子各需其特定的密度与温度，方能
发光：合起来便可对分子云作断层成像——外层包络在 CO 中显现，
致密核心则在 NH$_3$ 与 HCN 中显现。
\textbf{23.} 谱线是在宇宙背景的
衬托下被看到的；当气体温度趋近 \qty{2.7}{K} 时，其
激发与背景达到平衡，对比度随之
消失：这种方法无法读取更冷的分子云。
\textbf{24.} $\lambda/20 = \qty{130}{\micro m}$，相较于
光学工作的 $\qty{25}{nm}$，宽容度高出五千倍，这正是
十二米毫米波天线得以矗立于旷野，
而光学镜面须栖身圆顶之中。
\textbf{25.} 一个 $BJ(J+1)$ 形的阶梯，$B/h = \qty{57.6}{GHz}$，
在十开尔文下仍然活跃，于 \qty{2.6}{mm} 处读取，测得质量、
温度、湍流与转动，即冷宇宙的各项性质——而
其平坦的转动曲线是暗物质长久的例证。
\end{solution}
